% Generated draft: requires mathematical and source review.
\documentclass[11pt]{article}
\usepackage[margin=25mm]{geometry}
\usepackage{fontspec}
\setmainfont{Latin Modern Roman}
\usepackage{amsmath,amssymb,mathtools,unicode-math}
\setmathfont{Latin Modern Math}
\usepackage{enumitem,needspace,xurl,hyperref}
\setlength{\parindent}{0pt}
\setlength{\parskip}{5pt}
\setlength{\emergencystretch}{3em}
\newcommand{\N}{\mathbb{N}}
\newcommand{\Z}{\mathbb{Z}}
\newcommand{\Q}{\mathbb{Q}}
\newcommand{\R}{\mathbb{R}}
\newcommand{\C}{\mathbb{C}}
\newcommand{\F}{\mathbb{F}}
\newcommand{\A}{\mathbb{A}}
\newcommand{\PP}{\mathbb{P}}
\newcommand{\E}{\mathbb{E}}
\newcommand{\eps}{\varepsilon}
\newcommand{\dd}{\,\mathrm{d}}
\begin{document}
\section*{UnsolvedMath ID: 11100003}
\textbf{Source problem code:} catalogue entry
\subsection{Definitions and mathematical statement}
Hovey's Major problem 3 asks for a geometric meaning for elliptic cohomology. This is a programmatic request for geometric constructions or interpretations of the cohomology theories, not a single specified existence theorem or the assertion that no geometric interpretations are known.

The source requires the scope of the question to include Hopkins's viewpoint on elliptic spectra: elliptic cohomology comprises theories associated with different elliptic-curve/formal-group data, rather than only the Landweber--Ravenel--Stong theory based on the Jacobi quartic. In the customary schematic description, an elliptic spectrum relates a suitable periodic cohomology spectrum to an elliptic curve over its coefficient ring through an identification of their formal groups. The precise spectrum, periodicity, multiplicative structure, geometric cocycle model and compatibility requirements are not fixed in the excerpt.

Hovey expresses optimism that the problem may be solvable and states a strong referee preference for work engaging with the Hopkins viewpoint, criticizing treatment of the Jacobi-quartic example as the entire subject. This is attributed source commentary, not a mathematical hypothesis or an instruction governing this draft. He also points to Grojnowski's equivariant approach and Ando's work on power operations and loop-group representations, while acknowledging that other relevant work may be omitted.

The registry reports geometric interpretations through free loop spaces/Witten genus, derived elliptic moduli and supersymmetric field theories, but says that a fully satisfactory general interpretation, especially a Stolz--Teichner conjectural formulation, remains open. That conjecture is not stated here, and these different constructions are not silently equated or substituted for one uniquely specified target.
\subsection{Short English statement}
Seek geometric interpretations of elliptic cohomology in the broad elliptic-spectrum sense, encompassing the different theories rather than only the Jacobi-quartic example. The desired geometric model and compatibility axioms are unspecified.
\subsection{Sources}
\begin{itemize}
\item Local registry: \url{data/unsolvedmath/problems.json}, permanent UnsolvedMath ID $11100003$, exact source code \texttt{catalogue entry}, attributed to Mark Hovey.
\item \emph{Hovey's Algebraic Topology Problem List}, surviving mirror, Major problems item 3: \url{https://www-users.cse.umn.edu/~tlawson/hovey/}. Supplied access date: 2026-07-29; extraction: html-ordered-list; 65 top-level items in nine category pages.
\item M. J. Hopkins, ``Topological modular forms, the Witten genus, and the theorem of the cube,'' in \emph{Proceedings of the International Congress of Mathematicians, Vol. 1, 2 (Z\"urich, 1994)}, 554--565, Birkh\"auser, Basel, 1995; MR 97i:11043.
\item Grojnowski's approach to equivariant elliptic cohomology, described by the source as an unpublished preprint at its writing; no exact title or identifier supplied. M. Ando, ``Power operations in elliptic cohomology and representations of loop groups,'' \emph{Trans. Amer. Math. Soc.}, with source identifier printed \texttt{CNO CMP 1 637 129}; volume/year/pages are absent.
\item The registry mentions Stolz--Teichner and other geometric frameworks without exact further references or theorem statements.
\end{itemize}
\end{document}
